%---------------------------------------------------------------------
%====================================================================
\section{Summary of Fluid EFT}
%====================================================================
%-----------------------------------------------------------------
%
In this section we summarize the conventions and formulas of the fluid EFT for our convenience while drafting, regardless of what will actually appear in the main text.\notetag{This sec will be shortened at the end.}
First we present the fomrmulas valid at any order and then the perturbative expansion.


%====================================================================
\subsection{Effective larangian \& thermodynamic interpretation}\label{sec:fluid_eft_exact}
%====================================================================
%
The mapping from the spacetime manifold $\mathcal{M}^{(4)}$ to the internal fluid space $\mathcal{F}^{(3)}$ is given by four scalar fields\footnote{Note we use $\varphi$ for fluid coordinates, not to confuse them with Bardeen or Newtonian potentials.}  
{\small\begin{align}\label{eq:internal_coord_fluid}
    \varphi^{\alpha}:(\mathcal{M},g)\to \mathbb{R}\times \mathcal{F},\quad \alpha=0,1,2,3.
\end{align}}
Note we are free to choose the coordinates in the 3-dimensional internal space and the time coordinate $\varphi^0$ without using any gauge freedom of the spacetime manifold itself.
In particular, we might choose $\varphi^i(x)=x^i$ for $i=1,2,3$ and at some fixed initial time $x^0=\tau_{in}$.
We might also choose $\varphi^0(x)=\Xi(x^0)$ for some arbitrary function $\Xi$ at the background level, to be specified by demading suitable conservations, and then allow for perturbations $\pi^0(x^\mu)$ around it as in \eqref{eq:perturbed_fluid_coordinates}.
Most likely it will be convenient to choose $\Xi(x^0)=x^0$ so that the background value of $\varphi^0$ coincides with the time coordinate, but keep it general for now.\todotag{Fix when all settled}

The 4-velocity of the fluid is identified by requiring the spatial coordinates of the fluid be indeed comoving with it, that is $\partial_\mu \varphi^i=0$ for $i=1,2,3$.
After normalizing the 4-velocity we get
\begin{align}\label{eq:definition_4_velocity}
    u^\mu=& \frac{1}{\sqrt{-g}\, 3!\,b} \epsilon^{\mu \alpha \beta \gamma} \epsilon_{ijk} \, \partial_\alpha \varphi^i \partial_\beta \varphi^j \partial_\gamma \varphi^k\, \propto\,\,  \star(d\varphi^1\wedge d\varphi^2\wedge d \varphi^3), \quad u_\mu u^\mu=-1\,.
\end{align}
We can identify 2 adimensional quantities
\begin{align}\label{eq:definition_b_y}
    b=\sqrt{\det B}\quad \text{for}\quad
    B^{ij}= g^{\mu\nu} \partial_\mu \varphi^i \partial_\nu \varphi^j,\quad y= u^\mu \partial_\mu \varphi^0, \qquad [b]=[y]=[u^\mu]=E^0,
\end{align}
that respect Lorentz invariance in the spacetime manifold and the internal symmetries of the fluid, namely 
{\small\begin{align}\label{eq:symmetries_fluid_internal}
    &\varphi^i \mapsto \tilde{\varphi}^i\big(\{\varphi^j\}_j\big)\,\,\text{with}\,\, \det\Big(\frac{\partial \tilde{\varphi}^i}{\partial\varphi^j}\Big)=1\,\,\quad \text{(volume preserving 3d diffeo invariance)},\\
    &\varphi^0 \mapsto \varphi^0 + c\quad\qquad\quad\qquad\qquad\qquad\qquad \text{(shift symmetry in $\varphi^0$)}.
\end{align}}
Up to a \emph{dimensional} scale $\Lambda$ these variables are naturally identified with the number density $b\equiv n$ and the temperature $y\equiv T$ in the fluid EFT as explained in the next section.

The lagrangian is a generic function of the invariants
\begin{align}\label{eq:eft_action_exact}
    S=\int d^4x \sqrt{-g} F(b,y)\,.
\end{align}
The conserved Noether currents are the number density current (from volume-preserving diffeo invariance in the $\varphi^i$) and the entropy current\footnote{\label{fn:thermodynamic_interpretation}See next subsection for the jutification of the thermodynamic interpretation.} (from shift symmetry in $\varphi^0$)
\begin{align}\label{eq:noether_currents_nonbaro}
    \nabla_\mu J^\mu=0\quad\text{for}\qquad J_n^\mu = b\, u^\mu\,,\qquad J_s^\mu = F_y \,u^\mu\,.
\end{align}
Notably, combining the two equations above we see that the comoving entropy is conserved along the flow\footnote{Just use $\nabla_\mu(Qu^\mu)=\partial_\mu Q u^\mu+Q\nabla_\mu u^\mu$.}
\begin{align}\label{eq:comoving_entropy_conserved_along_flow}
    u^\mu \,\partial_\mu \sigma=0\qquad\text{for}\qquad \sigma:=\frac{s}{b}=\frac{F_y}{b}.
\end{align}
In a gauge comoving with the fluid $u^\mu\propto (1,000)$ and the above reads $\partial_0\sigma=0$.
In a general gauge, at the backround level we still have $\partial_0\bar{\sigma}=0$ and at perturbative order
\begin{align}\label{eq:entropy_conservation_perturbative}
    u^0\partial_0{\delta \sigma}-\delta u^j\, \partial_j\delta\sigma=0.
\end{align}
In particular we still have $\partial_0\delta\sigma=0$ at \emph{linear} perturbative order for a single fluid, but we anticipate this will not be the case for two fluids since relative entropy perturbations will become dynamical.

Varying the action wrt the metric, the energy-momentum tensor takes the perfect fluid form\footnote{Subscripts mean $X_b=\partial_b X_{\mid y}$ and $X_y=\partial_y X_{\mid b}$ i.e. main variables $(b,y)$, unless otherwise stated.}
\begin{align}\label{eq:em_tensor_nonbaro}
    T_{\mu\nu} = u_\mu u_\nu (yF_y - bF_b) + g_{\mu\nu} (F-bF_b) = u_\mu u_\nu(\rho+p) + g_{\mu\nu}\, p\,.
\end{align}
Notably the computation is exact to any order so that no anisotropic stress appear in the EFt of a single fluid.
Energy density, pressure and equation of state parameter are then identified as
\begin{align}\label{eq:thermo_rho_p_nonbaro}
    \rho=-F+yF_y,\quad p=F-bF_b,\quad w=\frac{F-bF_b}{yF_y-F}.
\end{align}
The entropy density $s$, comoving entropy $\sigma$, temperature $T$, chemical potential $\mu$ and number density $n$ are identified as\footnotemark[\getrefnumber{fn:thermodynamic_interpretation}]
\begin{align}\label{eq:thermo_other_variables_nonbaro}
    s=F_y,\quad \sigma:=\frac{s}{n}=\frac{F_y}{b},\quad T=y,\quad \mu:=F_b, \quad n=b, \quad F=-\rho+Ts=-\mathfrak{a}.
\end{align}
Notably this identifies $F(y,b)$ with minus the Helmoltz free-energy density {\small $\mathfrak{a}=\frac{U-TS}{V}$.}
Beware the identification $y=T$ is not with the usual temperature `in the Boltzmann sense', but rather more abstractly with `the variable conjugate to the entropy'.

It is convenient to have an expression for the relative varitions of $y$ versus $b$ at fixed comving entropy $\sigma$ or energy density $\rho$, since these \emph{define} adiabatic and non-adiabatic perturbations respectively,
\begin{align}\label{eq:variation_y_vs_b_fixed_entropy_energy_density}
    0&=d\sigma = \frac{F_{yy}}{b}dy - \frac{F_y-bF_{by}}{b^2}db\quad \Rightarrow\quad \frac{dy}{db}_{|\sigma}= \frac{F_y-bF_{by}}{bF_{yy}},\\[3pt]
    0&=d\rho=yF_{yy}dy+(yF_{yb}-F_b)db\quad \Rightarrow\quad \frac{dy}{db}_{|\rho}=\frac{F_b-yF_{yb}}{yF_{yy}}.
\end{align}
Beware \eqref{eq:comoving_entropy_conserved_along_flow} implies that $\sigma$ will be constant only at the background level or in a specific gauge.
Adiabatic sound speed $c_a^2$ and non-adiabatic pressure coefficient $\Gamma$ are then obtained plugging \eqref{eq:variation_y_vs_b_fixed_entropy_energy_density} into the definitions,\notetag{Write exp that makes manifest the limit $F_y=0$}
{\small\begin{align}\label{eq:adiabatic_cs_nonadiabatic_pressure}
\begin{aligned}
    &c_a^2=\frac{\partial p}{\partial \rho}_{|\sigma}=\frac{-y^2F_{yy}b^2F_{bb}+(yF_y-byF_{by})^2}{y^2F_{yy}(yF_y-bF_b)}, \\
    &\Gamma:=\frac{\partial p}{\partial \sigma}_{\mid\rho}= \frac{b^2F_{bb}y^2F_{yy}-(yF_y-byF_{by})(bF_b-byF_{yb})}{\tfrac{1}{by}y^2F_{yy}(yF_y-bF_b)}.
\end{aligned}
\end{align}}
Using $F=-\mathfrak{a}$ and $s=F_y=-\partial_T\mathfrak{a}$, we rewrite the adiabatic sound speed as
\begin{align}
    c_a^2:=
    &=\frac{-b^2F_{bb} F_{yy}+(F_y-bF_{yb})^2}{F_{yy}\,(yF_y-bF_b)} = \frac{n^2\,\partial_n^2\mathfrak{a}\,\partial_Ts + (s-n\partial_ns)^2}{\partial_Ts\,\,(\rho+p)}\,,
\end{align}
which makes it manifest that $c_a^2\geq0$ provided no ghosts $\rho+P>0$, the second principle of thermodynamics $\partial_Ts\geq 0$ and thermodynamic stability $\partial_n^2\mathfrak{a}\geq 0$ hold.

The adiabatic sound speed \eqref{eq:adiabatic_cs_nonadiabatic_pressure} always coincides with the coefficient $c_s^2$ of the gradient term in the quadratic action for phonons, after integrating out the non-dynamical entropy mode $\pi^0$,
\begin{align}
    c_a^2\equiv c_s^2\,,
\end{align}
which shows that absence of ghosts, thermodynamic stability and gradient instabilities are closely related.
This was expected: longitudinal phonons, i.e. sound waves, do propagate with the speed of sound!
This is proved in \eqref{eq:cs_equiv_ca_flat_spacetime} for flat spacetime and in \eqref{eq:quadrati_action_final_integrated_out_entropy_field} for curved spacetime.


Beware that, whenever $p$ is not a function of $\rho$ only, we have\questiontag{Polish \& shorten here below}
\begin{align}
    \underbrace{c_s^2}_{\text{grad coef}}\equiv c_a^2:=\frac{\partial p}{\partial \rho}_{|\sigma}\neq \frac{\dot{p}}{\dot{\rho}}=:c_g^2.
\end{align}
Indeed for two independent thermodynamic variables, the ratio of time variations of pressure to density corresponds to taking derivatives along a \emph{curved} trajectory in the thermodynamic plane, i.e. directional derivatives $\tfrac{\dot{p}}{\dot{\rho}}=\tfrac{\partial P}{\partial \hat{v}(t)}$ along a time-dependent direction $\hat{v}(t)$ ultimately depending on the background evolution of the fluid.
In the general case, this procedure has no fixed relation with derivatives with a well-defined thermodynamic variable kept fixed.

On the other hand, in the barotropic case $p=p(\rho)$ the expression $dp/d\rho$ is unambiguous since variations of $p$ and $\rho$ are proportional with the same coefficient $\tfrac{dp}{d\rho}$ for any direction  thermodynamic plane,
\begin{align}\label{eq:barotropic_speed_of_sound}
    \frac{dp}{d\rho}=\frac{\tfrac{dp}{d\hat{x}}}{\tfrac{d\rho}{d\hat{x}}}=\frac{\tfrac{dp}{d\hat{z}}\cancel{\tfrac{d\hat{z}}{d\hat{x}}}}{\tfrac{d\rho}{d\hat{z}}\cancel{\tfrac{d\hat{z}}{d\hat{x}}}}=\frac{\tfrac{dp}{d\hat{z}}}{\tfrac{d\rho}{d\hat{z}}} = \frac{\dot{p}}{\dot{\rho}}=w-\frac{\dot{w}\bar{\rho}}{\dot{\bar{\rho}}}\quad \forall  \text{ directions } \hat{x},\hat{z}, \quad \text{and}\qquad c_a^2=c_s^2=c_g^2.
\end{align}
In particular, in the barotropic case adiabatic sound speed, gradient coefficient and ratio of time variations of pressure to density all coincide.

Moreover, if the relative variation of the EoS parameter is small compated to the Hubble rate
\begin{align}
    \frac{\tfrac{\dot{w}}{w}}{\tfrac{\dot{\bar{\rho}}}{\bar{\rho}}}\simeq \frac{\tfrac{\dot{w}}{w}}{-3\mathcal{H}(1+w)}\ll 1 \quad \Rightarrow \quad c_s^2\simeq w,
\end{align}
the EoS parameter and the gradient coefficient esentially coincide, so that accelerated expansion $w<-1/3$ unavoidably leads to gradient instabilities $c_s^2<0$.
That is, a barotropic fluid cannot support long-lasting accelerated expansion without instabilities!

For the record, comparing our formalism to Wayne Hu's \cite[equation (3)]{Hu:1998kj} we find
\begin{align}
    \delta p\overset{\text{us}}{=} \underbrace{\tfrac{\partial p}{\partial \rho}_{|\sigma}}_{=:c_a^2}\delta\rho + \underbrace{\tfrac{\partial p}{\partial \sigma}_{|\rho}}_{=:\Gamma}\delta\sigma \overset{\text{Hu}}{=} \underbrace{\tfrac{\dot{p}}{\dot{\rho}}}_{=:c_g^2} \delta \rho + \Gamma_{\text{Hu}} p
    \quad\Rightarrow\quad \Gamma_{\text{Hu}} p = \Gamma \delta\sigma + (c_a^2-c_g^2)\delta\rho.
\end{align}
Wayne's definition of $\Gamma_{Hu}$ is thus a mixture of energy density and entropy perturbations.
Recalling that entropy perturbations and velocity divergence are proportional $\delta\sigma \propto \theta$ and dividing the above equation by the energy density, we find
\begin{align}
    w\Gamma_{\text{Hu}} = \Gamma \underbrace{\tfrac{\delta\sigma}{\rho}}_{\propto\, \theta}+ (c_a^2-c_g^2)\,\delta,
\end{align}
which is precisely the expression given in Hu's \cite[equation (3)]{Hu:1998kj} in the rest frame of the fluid where $\theta=0$.


Finally we remark that a fluid is \emph{barotropic} if $p=p(\rho)$ only.
The case $F=F(b)$ is thus a special case of barotropic fluid, but not the most general one!
Indeed we can keep the barotropic condition $p=p(\rho)$ and still allow for  $F=F(b,y)$.
The simplest example is a constant EoS parameter, but more generally any $w=w(\rho)=w(-F+yF_y)$ works.

When $F=F(b)$, we can still define $y=\partial_\mu\varphi^0u^\mu$, but since $F$ is independent of $y$, the above formulas reduce to
{\small\begin{align}\label{eq:thermo_variables_barotropic}
\begin{aligned}
    &\rho=-F,\quad p=F-bF_b,\quad w=\frac{bF_b-F}{F},\quad n=b,\quad \mu=F_b,\\[3pt]
    & \frac{d p}{d \rho}=c_a^2=\frac{b^2F_{bb}}{bF_b}= w+w_b \frac{\rho}{\rho_b},\qquad T=y,\quad s=\sigma=0,\quad \Gamma=\frac{\partial p}{\partial \sigma}_{|\rho}=0.
\end{aligned}\end{align}}
As discussed above for the general barotropic case, there is now only one independent thermodynamic variable, which can be any of $b$, $\rho$ etc, and 
There is no ambiguity in the variables to keep fixed when computing $dp/d\rho$ and \eqref{eq:barotropic_speed_of_sound} holds, which now further reduces to
\begin{align}
    c_s^2\overset{\substack{\text{shown}\\ \text{in general}}}{=}
c_a^2:=\frac{dp}{d\rho}\overset{F_y=0}{=}\frac{\tfrac{dp}{db}\cancel{db}}{\tfrac{d\rho}{db}\cancel{db}}= \frac{b^2F_{bb}}{bF_b}=\frac{\tfrac{dp}{db}\cancel{\tfrac{db}{dt}dt}}{\tfrac{d\rho}{db}\cancel{\tfrac{db}{dt}dt}}=
    \frac{\dot{\bar{p}}}{\dot{\bar{\rho}}}=:c_g^2= w-\frac{\dot{w}\bar{\rho}}{\dot{\bar{\rho}}}.
\end{align}



%-----------------------------------------------------------------------
\subsubsection{Justification of the thermodynamic interpretation}
%----------------------------------------------------------------------
In this section we justify the thermodynamic interpretation of the fluid EFT variables and formulas.
As argued above, the identification of the 4-velocity \eqref{eq:definition_4_velocity} follows unavoidably from the requirement the spatial coordinates $\varphi^i$ be indeed comoving with the fluid.

The next point is the unequivocal identification of $b$ with the number density $n$.
Indeed, in order for the mapping \eqref{eq:internal_coord_fluid} from spacetime to the internal fluid space to make sense, it is necessary that the restriction $\varphi^{I}:\Sigma_\tau^{(3)}\to \mathcal{F}^{(3)}$ is a diffeomorphism for any spacelike hypersurface $\Sigma_\tau^{(3)}\subset \mathcal{M}^{(4)}$.
The fluid coordinates $\{\varphi^i\}_{i=1,2,3}$ can thus be used as coordinates for the hypersurface $\Sigma_\tau^{(3)}$, and $g^{\mu\nu}\partial_\nu\varphi^i$ are 3 independent vectors spanning its tangent space.
In these coordinates, the metric $g_{|\Sigma}$ induced on spatial slices reads
\begin{align}
    (g^\varphi_{\Sigma})_{ij} = \partial_\mu \varphi^i\partial_\nu \varphi^j g^{\mu\nu}\equiv B^{ij}.
\end{align}
The volume factor $\sqrt{g^\varphi_\Sigma}=\sqrt{\det B}=b$, i.e. the jacobian of the transformation from spatial coordinates to internal fluid coordinates, is then precisely the density of the fluid elements wrt the original spacetime coordinates $b\equiv n$.

To complete the thermodynamic interpretation we need some final imput, which comes from matching the EFT energy-momentum tensor to that of a perfect fluid \eqref{eq:em_tensor_nonbaro}, thereby identifying the energy density and pressure \eqref{eq:thermo_rho_p_nonbaro}.
Finally we simply impose the first law of thermodynamics
\begin{align}\label{eq:first_law_thermo}
    T s+ \mu n = \rho+p \overset{!}{=} yF_y-bF_b.
\end{align}
Since we have already shown $n=b$, we are lead to identify $\mu=-F_b$, and $s=F_y$ and $y=T$.
This is consistent with all the thermodynamical relations, for example imposing the variation of the entropy density vanish we find back $\mu$ from its thermodynamic definition
\begin{align}\label{eq:chemical_potential_and_entropy_consistency}
    \begin{array}{c}
     0\overset{!}{=}ds = F_{yb} db + F_{yy}dy\\[3pt]
    \displaystyle\Rightarrow\,\, \frac{dy}{db}_{|s}=-\frac{F_{yb}}{F_{yy}}
    \end{array}
    \quad\Rightarrow\quad \mu:=\frac{\partial\rho}{\partial n}_{|s}= \underbrace{\frac{\partial\rho}{\partial b}_{|y}}_{=-F_b+yF_{by}}+ \underbrace{\frac{\partial\rho}{\partial y}_{|b}}_{=yF_{yy}}\,\,\frac{dy}{db}_{|s}=-F_b.
\end{align}
In fact, after enforcing $n=b$ and the expressions for $\rho$ and $p$, equating $s=F_y,\,T=y$ is the \emph{only} consistent identification.
Indeed taking $n=b$ and $y=s$ predicts wrong conjugate variables and fails the 1st law of thermodynamics \eqref{eq:first_law_thermo} since
\begin{align}\begin{aligned}
    &n=b,\,\,s= y \,\,\Rightarrow\,\, \begin{cases}
        \mu:=\frac{\partial\rho}{\partial n}_{|s}= \frac{\partial\rho}{\partial b}_{|y}=-F_b+yF_{by}\\
         T:=\frac{\partial\rho}{\partial s}_{|n}= \frac{\partial\rho}{\partial y}_{|b}=yF_{yy}
    \end{cases}\\[4pt]
    &\,\,\Rightarrow\,\, \mu n +T s= -bF_b+byF_{by}+y^2F_{yy}\neq yF_y-bF_b=\rho+p.
\end{aligned}\end{align}
Finally we stress the identification $y=T$ is not with the familiar temperature `in the Boltzmann sense', but rather with the a `generalized temperature' to be interpreted as the variable conjugate to the generalized entropy density $s=F_y$.
In fact even this notion of entropy density has wrong dimension $[F_y]=E^4$ and should be rescaled in order to identify it familiar one.
In turn, the proportionality of $y$ versus the scale factor $a$ is in general different from the one of the Boltzmann temperature as discussed after \eqref{eq:bkg_scaling_y_coordinate}.


%=============================================================================
\subsection{Background evolution in the fluid EFT}\label{sec:bkg_evolution_fluid_EFT}
%==============================================================================
%
We consider first the conservation of the Noether currents.
Since $u^\mu=a^{-1}\delta^\mu_0$ and $\sqrt{g}=a^{4}$ at the background level, we have
\begin{subequations}\label{eq:conservation_noether_current_background_eft}
\begin{align}
    &0=\nabla_\mu J_n^\mu \frac{1}{\sqrt{g}} \partial_\mu\Big(\sqrt{-g} b u^\mu\Big) \revpropto \frac{d}{d\tau}(a^3 b)\quad \Rightarrow \quad b = b_0 a^{-3},\label{eq:conservation_noether_current_background_eft_number}\\[3pt]
    &0=\nabla_\mu J_s^\mu =\frac{1}{\sqrt{g}} \partial_\mu\Big(\sqrt{-g} F_y u^\mu\Big) \revpropto \frac{d}{d\tau}(a^3 F_y)\quad \Rightarrow \quad F_y\propto a^{-3}. \label{eq:conservation_noether_current_background_eft_entropy}
\end{align}
\end{subequations}
Using \eqref{eq:conservation_noether_current_background_eft_number}, the second equation \eqref {eq:conservation_noether_current_background_eft_entropy} can be equivalently rewritten as an evolution equation for $y$,
\begin{align}\label{eq:eolvution_y_background}
    \dot{y}=\underbrace{\dot{b}}_{=-3\mathcal{H}b}\,\frac{F_y-bF_{by}}{bF_{yy}} =\dot{b}\, \frac{dy}{db}_{|\sigma}\,.
\end{align}
The second equality follows from \eqref{eq:variation_y_vs_b_fixed_entropy_energy_density} and implies that the evolution of $y$ is adiabatic at the background level.
This is no surprise since \eqref{eq:comoving_entropy_conserved_along_flow} alredy showed that $\dot{\bar{\sigma}}=0$ at the background level.
Finally, from the very definition of $y=u^\mu\partial_\mu\varphi^0$ we have the background evolution
\begin{align}\label{eq:evolution_y_background_from_definition}
    \bar{y}=\partial_\mu\bar{\varphi}^0 \bar{u}^\mu = \frac{1}{a}\mathcal{H}\frac{d\Xi}{d\log\!a},
\end{align}
with $y\propto a^{-1}$ in the natural case where $\Xi(\tau)=\tau$.


{\color{gray}
----------------------------------------------------------------------------------------\todotag{Ignore for the moment}

Crucially \eqref{eq:eolvution_y_background} is equivalent to the characteristic ODE \eqref{eq:characteristics_eos_pde_w_depending_b}-\eqref{eq:characteristics_eos_pde_solution_w_depending_b} of the equation of state with $w=w(b)$.
Evolution in time thus coincides with evolution along characteristics of the equation of state \eqref{eq:eos_pde_w_depending_b}.
This was expected since $a\propto\, b^{-1/3}$ from \eqref{eq:conservation_noether_current_background_eft_number} implies we can use either variable to parametrize time.
Indeed, from the general solution \eqref{eq:eos_pde_solution_w_depending_b} of the EoS we compute
\begin{align}
    F(y,b)=f\Big(y\, e^{-\int_{b_0}^bd\log b\, w(b)}\Big)\, e^{\int_{b_0}^bd\log b\, \big(1+w(b)\big)} \quad \Rightarrow\quad  \frac{F_y-bF_{by}}{bF_{yy}} \equiv w(b)\frac{y}{b}.
\end{align}
Plugging this into \eqref{eq:eolvution_y_background} we get precisely the characteristic ODE \eqref{eq:characteristics_eos_pde_w_depending_b}, now written in terms of both the equivalent variables $b$ or $a$,
{\small\begin{align}\label{eq:evolution_y_background_equation_final}
\begin{aligned}
    \frac{\dot{y}}{y}=w(b)\frac{\dot{b}}{b}=-3w(a)\mathcal{H}.
\end{aligned}
\end{align}}
The solution is of course the same as \eqref{eq:characteristics_eos_pde_solution_w_depending_b}, again rewritten in terms of $a$ or $b$ respectively,
{\small\begin{align}\label{eq:evolution_y_background_solution}
\begin{aligned}
    y
    &=y_0\,\,\exp\Big(\!\!-\!\!3\!\!\int^a_{a_0}\!\!\!\!\!d\!\log(a)\,\, w(a)\Big)\\[2.5pt]
    &=y_0\,\exp\Big(\!\!-\!\!\!\int_{b_0}^b\!\!\!\!\!d\!\log\big(\tfrac{b}{b_0}\big)\,w(b)\Big)
    =y_0\,\underbrace{\left(\tfrac{b}{b_0}\right)^{w(a)}}_{=a^{-3w(a)}}\,\exp\Big(\!\!-\!\!\!\int_{a_0}^a\!\!\!\!\!da\, \tfrac{dw}{da}\, \underbrace{\log\big(\tfrac{b}{b_0}\big)}_{=-3\log a}\Big).
\end{aligned}
\end{align}}
We can now specify the background value $\varphi^0(x^\mu)=\Xi(x^0)+\pi^0(x^\mu)$ for the time component of the fluid internal coordinates \eqref{eq:perturbed_fluid_coordinates}.
Enforcing it be consistent with \eqref{eq:evolution_y_background_solution} we find
\begin{align}\label{eq:bkg_scaling_y_coordinate}
\begin{aligned}
    &y=u^\mu\partial_\mu\varphi^0 \overset{\text{bkg}}{=}\tfrac{1}{a}\mathcal{H}\frac{d\Xi}{d\log\!a}\propto\, \exp\Big(\!\!-\!\!3\!\!\int^a_{a_0}\!\!\!\!\!d\!\log(a)\,\, w(a)\Big)\\[2.5pt]
    &\quad\Rightarrow\quad \Xi(a)=\int^a_{a_0}\!\!\!d\log\!a\, \frac{a}{\mathcal{H}}e^{-3\int^a_{a_0}\!d\!\log(a)\,\, w(a)}\,,
\end{aligned}\end{align}
up to a nonphysical constant shift in time $\Xi\mapsto \Xi +C$.

This clarifies the identification $y=T$ is not with the familiar temperature `in the Boltzmann sense', but rather with the a `generalized temperature' to be interpreted as the variable conjugate to the generalized entropy density $s=F_y$.
In fact even this notion of entropy density has wrong dimension $[F_y]=E^4$ and should be rescaled in order to identify it familiar one.
In turn, beware the proportionality of $y$ versus the scale factor $a$ is generally different from the one of the Boltzmann temperature.
It accidentally coincides for radiation $w=\tfrac13,\, y\propto \tfrac{1}{a}$ and cold matter $w=0,\,y\propto\, a^0$, but is different for any species in between, e.g. the familiar scaling for non-relativistic matter $y\propto\, a^{-2}$ would require superluminal speed of sound $w=\tfrac23$.
-------------------------------------------------------------------------------
}

The continuity equation for the background is automatically implied by the conservation of Noether currents, and yields no further constraint or information.
Indeed, using expressions \eqref{eq:thermo_rho_p_nonbaro} for $\rho$ and $p$, the continuity equation \eqref{eq:background_continuity} reads
\begin{align}\label{eq:backgroud_continuity_eft}
    \frac{d}{d\tau}(-F+yF_y)+3\mathcal{H}(yF_y-bF_b)=0.
\end{align}
Expanding $\dot{F}=F_b\dot{b}+F_y\dot{y}$ and $\dot{F}_y=F_{yb}\dot{b}+F_{yy}\dot{y}$, one verifies the above is indeed implied by \eqref{eq:conservation_noether_current_background_eft}.
In the barotropic case \eqref{eq:thermo_variables_barotropic} it reduces to the familiar $a^{-3}$ dilution of the number density $b$.
For the record, we remark there is no advantage in plugging the explicit solution \eqref{eq:constant_eos_w_solution_pde} or \eqref{eq:eos_pde_solution_w_depending_b} into the background continuity equation.
It just gives a fairly messy expression, whose information is already contained and conveniently split in the conservation of the two Noether currents.

{\color{gray}-------------------------------------------------------\todotag{Ignore for the moment}

Having enforced the desired EoS parameter $w(a)$ by writing $F$ as \eqref{eq:eos_pde_solution_w_depending_b}, and having demanded the conservation of Noether currents in \eqref{eq:conservation_noether_current_background_eft} and thus the continuity equation, we finally impose the last background equation (1st Friedmann).
Inserting the expression \eqref{eq:eos_pde_solution_w_depending_b} into $\rho=-F+yF_y$ and then plugging the time evolution \eqref{eq:evolution_y_background_solution} of $y$ and \eqref{eq:conservation_noether_current_background_eft_number} of $b$, we get
\begin{align}
\begin{aligned}
    \rho &= yb\, f^\prime\Big(y\, e^{-\int_{b_0}^bd\log\!b\, w}\Big)-e^{-\int_{b_0}^bd\log\!b\, (1+w)}\,f\Big(y\, e^{-\int_{b_0}^bd\log \!b\, w}\Big) \\
    &= e^{-\int_{b_0}^bd\log\!b\, \big(1+w(b)\big)}\underbrace{\Big[y_0f^\prime(y_0)-f(y_0)\Big]}_{=\rho_0=\text{const}} = \rho_0 \, a^{-3}\, e^{-3\int_{a_0}^a\!\!d\log\!a\, w(a)}.
\end{aligned}
\end{align}
The energy density is expressed as a function of the scale factor $a$, also through the form of $w(a)$, and the initial condition $\rho_0=y_0f^\prime(y_0)-f(y_0)$.

The First Friedmann equation then reads
\begin{align}
    &\mathcal{H}^2 = \frac{8\pi G}{3} a^2 \rho = \frac{8\pi G}{3} a^2 \rho_0 \, a^{-3}\, e^{-3\int_{a_0}^a\!\!d\log\!a\, w(a)}.
\end{align}
For a given functional form $w(a)$ we can now solve the above to get the time evolution of the scale factor $a(\tau)$.
The case of familiar species is obtained for constant $w=-1,0,1/3$.
---------------------------------------------------------------------}




%==============================================================
\subsection{Perturbative expansion \& matching phonons to fluid perturbations}
%============================================================
%
The formulas of Section \ref{sec:fluid_eft_exact} are exact to any order in the fluid coordinates.
We now expand to quadratic order around unperturbed\footnote{Beware this choice of fluid coordinates does not use any gauge freedom of the spacetime manifold!} fluid $\bar{\varphi}^\mu$ and spacetime $x^\mu$ coordinates.
The general perturbed FLRW metric is written wrt conformal time $x^0=\tau$ as
{\small \begin{align}\label{eq:general_perturbed_FLRW_metric}
\begin{aligned}
    &g_{\mu\nu}=a^2\Big(\eta_{\mu\nu}+h_{\mu\nu}\Big)\,,\quad g^{\mu\nu}=a^{-2}\Big(\eta^{\mu\nu}-h^{\mu\nu}\Big)\quad \text{with} \quad h^{\mu\nu}=\eta^{\mu\alpha}\eta^{\nu\beta}h_{\alpha\beta},\\
    &\quad \sqrt{|g|}= a^4\big(1+\tfrac12 h^{\mu\nu}\eta_{\mu\nu} +O(h^2)\big)\ = a^4\Big[1+\tfrac12 (h^{ii}-h^{00}) +O(h^2)\Big]\,.
\end{aligned}
\end{align}}
The fluid variables are written for simplicity with rescaling factors $\gamma$ and $\lambda$ that will be used to reabsorb units of measures,
{\small\begin{align}\label{eq:perturbed_fluid_coordinates}
    \varphi^0=\gamma\Big(\tau+\pi^0\Big), \quad \varphi^i=\lambda\Big(x^i+\pi^i\Big)\,, \quad [\varphi^\alpha]=[x^\alpha]=[\pi^\alpha]=E^{-1},\quad [\gamma]=[\lambda]=E^0\,.
\end{align}}
Since internal fluid coordinates are independent of spacetime coordinates, under a gauge transformation $x^\mu \mapsto x^\mu +\delta x^\mu$ with $\delta x^\mu=(T,\,\partial^i L+\hat{L}^i)$, the fluid perturbations must tranform as
{\small\begin{align}\label{eq:gauge_transformation_fluid_perturbations}
\begin{aligned}
    &\pi^0 \mapsto \pi^0 - T,\qquad \pi^i \mapsto \pi^i - \partial^i L - \hat{L}^i.
\end{aligned}\end{align}}
Using the gauge transformation of the metric \eqref{eq:gauge_transformation_metric}, the relevant combinations change as
{\small
\begin{align}\label{eq:gauge_transformation_fluid_and_metric_combinations}
\begin{aligned}
    &\dot{\pi}^0+\tfrac12 h^{00} \mapsto \dot{\pi}^0+\tfrac12 h^{00} +\mathcal{H} T, \\
    &\partial_j\pi^j -\tfrac12 h^{ii} \mapsto \partial_j\pi^j -\tfrac12 h^{ii} -3\mathcal{H} T, \\
    &\dot{\pi}^j+h^{0j} \mapsto \dot{\pi}^j+h^{0j} -\partial^j T,
\end{aligned}\end{align}}
and we can combine these two together or with $\mathcal{H}\frac{\delta \rho}{\dot{\rho}}$ from \eqref{eq:gauge_transformation_em_tensor} to get gauge invariant combinations.\todotag{Write em \& make compare to Bardeen ones}

The invariants $b$ and $y$ are expanded as follows
{\small\begin{align}\label{eq:expansions_b_y}
\begin{aligned}
    &b= \bar{b}(1+\delta b),\quad \bar{b}= \frac{\lambda^3}{a^3},\quad
    \delta b^{(1)} = -\tfrac12 h^{ii} +(\partial_j \pi^j),\\[3pt]
    &\delta b^{(2)} = - h^{0 i} \dot{\pi}^i - \tfrac12  h^{ii} (\partial_j \pi^j) - \tfrac12 (\dot{\pi}^i)^2 + \tfrac12 (\partial_j \pi^j)^2 -\tfrac12 \partial_i\pi^j\partial_j\pi^i\,,
    \\[3pt]
    &y= \bar{y}(1+\delta y),\quad \bar{y}= \frac{\gamma}{a},\quad \delta y^{(1)} = \tfrac12 h^{00}+ \dot{\pi}^0,\\[3pt]
    &\delta y^{(2)} = \tfrac12 h^{00}\dot{\pi}^0+ h^{0 i} \dot{\pi}^i  + \tfrac12 (\dot{\pi}^i)^2 -\partial_j\pi^0\dot{\pi}^j\,.
\end{aligned}\end{align}}
The 4-velocity $u^\mu$ is also expanded, highlighting both the covariant form and the explicit splitting in temporal and spatial part,
{\small\begin{align}\label{eq:expansion_u_4_velocity}
\begin{aligned}
    au^\mu&=\delta_0^\mu \bigg(1 \!+\! \Big[\tfrac12 h^{00} -(\partial_j \pi^j)\Big]\! +\! \Big[ h^{0 i} \dot{\pi}^i -\tfrac12 h^{00} (\partial_j \pi^j) + \tfrac12 (\dot{\pi}^i)^2 + \tfrac12 (\partial_j \pi^j)^2 +\tfrac12 \partial_i\pi^j\partial_j\pi^i\Big]\bigg)
    \\
    &\quad\qquad+\tfrac{1}{2}\left[1 + \tfrac12 h^{00}-(\partial_\ell\pi^\ell)\right]
     \epsilon_{i j k}\,\epsilon^{\mu \alpha j k}\,\partial_\alpha\pi^i 
     +\frac{1}{2} \epsilon_{ijk}\, \epsilon^{\mu i \beta\gamma}\partial_\beta\pi^j\,\partial_\gamma\pi^k
    + O(\pi^3, h^2)\\[3pt]
    &=\Big(1 + \tfrac12 h^{00} +h^{0i}\dot{\pi}^i+ \tfrac12 (\dot{\pi}^i)^2,\,\,- \dot{\pi}^\ell + \dot{\pi}^m \partial_m\pi^\ell - \tfrac12 h^{00}\dot{\pi}^\ell\Big)+ O(\pi^3, h^2).
\end{aligned}\end{align}}
%
We also introduce the phonon potential to rewrite longitudinal (scalar) phonons
\begin{align}\label{eq:phonon_potential}
    \pi_L=\nabla^{-2}\partial_\ell\pi^\ell, \quad \pi^\ell_{\text{||}}=\partial_\ell \pi_L,\quad  [\pi_L]=E^{-2}.
\end{align}
Matching the velocity expansion \eqref{eq:expansion_u_4_velocity} in the fluid EFT to the general expression \eqref{eq:velocity_decomposition} in pertrubed FLRW spacetime $au^\mu =(1+.., v^\ell)$, we identify the peculiar velocity and the velocity potential/divergence in terms of phonons as
\begin{align}\label{eq:peculiar_velocity_fluid_vs_phonons}
    v^\ell = \frac{1}{2}\epsilon_{i j k}\,\epsilon^{\ell \alpha j k}\,\partial_\alpha\pi^i = -\dot{\pi}^\ell\quad \Rightarrow \quad \theta=-\partial_\ell\dot{\pi}^\ell= -\nabla^2\dot{\pi}_L,\,\, \,\,v = -\dot{\pi}_L\,.
\end{align}

Finally we recast the fluctuations in density $\rho=-F+yF_y$, pressure $p=F-bF_b$ and comoving entropy $\sigma=\tfrac{F_y}{b}$ as
\begin{align}\label{eq:rho_p_fluctuations_via_b_y}
\begin{aligned}
    \delta\rho &= y^2F_{yy}\delta y + (-bF_b+byF_{by})\delta b
    \\
    &= (y^2F_{yy})\left(\tfrac12 h^{00}+ \dot{\pi}^0\right) + (-bF_b+byF_{by})\left(-\tfrac12 h^{ii} +(\partial_j \pi^j)\right) + O(2)\,,
    \\[3pt]
    \delta p &= (yF_y-byF_{by})\delta y - (b^2 F_{bb})\delta b
    \\
    &= (yF_y-byF_{by})\left(\tfrac12 h^{00}+ \dot{\pi}^0\right) - (b^2 F_{bb})\left(-\tfrac12 h^{ii} +(\partial_j \pi^j)\right)+ O(2)\,,
    \\[3pt]
    \delta\sigma &= \tfrac{yF_{yy}}{b}\delta y + \tfrac{bF_{yb}-F_y}{b} \delta b
    \\
    &= \tfrac{y^2F_{yy}}{yb}\left(\tfrac12 h^{00}+ \dot{\pi}^0\right) +\tfrac{ybF_{yb}-yF_y}{yb}\left(-\tfrac12 h^{ii} +(\partial_j \pi^j)\right)+ O(2)\,.
\end{aligned}
\end{align}
We now express pressure fluctuations in terms of energy density and entropy ones.
Using the flat \eqref{eq:integrate_out_entropy_flat_spacetime} or curved spacetime expression \eqref{eq:integrate_out_entropy_curved_spactime} to integrate out the combination $\left(\tfrac12 h^{00}+ \dot{\pi}^0\right)$ of entropy field and metric perturbations, and recalling the expressions \eqref{eq:adiabatic_cs_nonadiabatic_pressure} for $c_a^2$ and $\Gamma$, after some math we find
\begin{align}\label{eq:rho_p_sigma_fluctuations_vs_phonon_entropy_field}
\begin{aligned}
    \delta\rho &= (yF_y-bF_b)(\partial_j\pi^j-\tfrac12 h^{ii})+a^{-4}E(\mathbf{x}) = (\bar{\rho}+\bar{p})(\partial_j\pi^j)\,+by\,\delta\sigma \,,\\[2.5pt]
    \delta\sigma &= \frac{a^{-4}}{by}E(\mathbf{x})\,,\\[2.5pt]
    \delta p &= c_a^2\, \delta\rho + \Gamma \, \delta\sigma\,.
\end{aligned}
\end{align}
As expected, phonons $\partial_j\pi^j$ only source density fluctuations, while the entropy field $\pi^0$ sources both density and entropy fluctuations \todoAC{This is for flat spacetime, write general form in curved spacetime}
\begin{align}\label{eq:entropy_field_vs_density_entropy_fluctuations}
\begin{aligned}
    \dot{\pi}^0&= \frac{1}{yF_y-byF_{by}}\bigg[\delta\rho \Big(c_a^2+\tfrac{b^2F_{bb}}{yF_y-bF_b}\Big)+\delta\sigma\,\Big(\Gamma-\tfrac{b^2F_{bb}}{\tfrac{1}{by}(yF_y-bF_b)}\Big)\bigg]\\[2.5pt]
    &=\tfrac{1}{y\partial_yp}\bigg[\delta\rho \Big(c_a^2-\tfrac{b\partial_bp}{\rho+p}\Big)+\delta\sigma\,\Big(\Gamma+by\tfrac{b\partial_bp}{\rho+p}\Big)\bigg]\,.
\end{aligned}
\end{align}






%------------------------------------------------------------------------
\subsubsection{Continuity \& Euler equation in the fluid EFT}\todotag{Finish section}
%---------------------------------------------------------------------
%
We now write the perturbed conservation equations $\nabla_\mu T^{\mu}_\nu=0$ in the fluid EFT.
We use the decomposition of the energy-momentum tensor \eqref{eq:energy_momentum_tensor_decomposition} and the EFT expansions \eqref{eq:rho_p_fluctuations_via_b_y} and \eqref{eq:peculiar_velocity_fluid_vs_phonons}, and then plug everything into the chosen gauge or even the general metric \eqref{eq:general_perturbed_FLRW_metric}.

\paragraph{Newtonian gauge}
For simplicity we start in Newtonian gauge \eqref{eq:newtonian_gauge_metric} and simply plug the relevant expressions into \eqref{eq:perturbed_continuity_drho}-\eqref{eq:euler_divergence_form}
\begin{align}
    ...
\end{align}



%==========================================================================
%\subsection{Quadratic action for phonons}
%==========================================================================
%
\input{temp_appendix/quadratic_action_NEW.tex   }


\newpage
%================================================================================
\subsection{Related computations (to be removed at the end)}
%================================================================================


%-------------------------------------------------------------------------
\subsubsection{Thermo formulas adiabatic sound speed \& non-adiabatic pressure}
%--------------------------------------------------------------------------

The entropy per particle is $\sigma=F_y/b$, imposing it be preserved gives
\begin{align}
    0=d\sigma = \frac{F_{yy}}{b}dy+\Big(\frac{bF_{by}-F_y}{b^2}\Big)db\quad \Rightarrow\quad \frac{dy}{db}_{\mid\sigma}= \frac{F_y-b F_{yb}}{bF_{yy}}=\frac{y}{b}\frac{p_y}{\rho_y}.
\end{align}
The adiabatic sound speed is then by definition\notetag{Polish redundant formulas \& make barotropic limit manifest}
\begin{align}
    c_b^2:=\frac{dp}{d\rho}_{\mid\sigma}&=\frac{p_b + p_y\, \tfrac{dy}{db}_{\mid\sigma}}{\rho_b + \rho_y\, \tfrac{dy}{db}_{\mid\sigma}} 
    =\frac{-b^2F_{bb} F_{yy}+(F_y-bF_{yb})^2}{F_{yy}\,(yF_y-bF_b)}\\[5pt]
    &= \frac{b p_b\, y\rho_y+(yp_y)^2}{b\rho_b\,y\rho_y+y\rho_y\,yp_y} = \frac{b\,p_b \,y\rho_y + (y\,p_y)^2}{y \rho_y (\rho+p)}\,.
\end{align}
The rewritings in the second line might come handy when trying to show positivity or other relations.
Using $F=-\mathfrak{a}$ and $s=F_y=-\partial_T\mathfrak{a}$, we rewrite the adiabatic sound speed as
\begin{align}
    c_b^2:=
    &=\frac{-b^2F_{bb} F_{yy}+(F_y-bF_{yb})^2}{F_{yy}\,(yF_y-bF_b)} = \frac{n^2\,\partial_n^2\mathfrak{a}\,\partial_Ts + (s-n\partial_ns)^2}{\partial_Ts\,\,(\rho+p)}\,,
\end{align}
which makes it manifest that $c_b^2\geq0$ provided no ghosts $\rho+P>0$, the second principle of thermodynamics $\partial_Ts\geq 0$ and thermodynamic stability $\partial_n^2\mathfrak{a}\geq 0$ hold.
The equivalence $c_a^2\equiv c_s^2$ is proved in \eqref{eq:cs_equiv_ca_flat_spacetime} for flat spacetime and in \eqref{eq:quadrati_action_final_integrated_out_entropy_field} for curved spacetime.

Similarly we compute the nonadiabatic pressure coefficient $\frac{\partial P}{\partial \sigma}_{|\rho}$.
We proceed as before enforcing 
\begin{equation}
    0=d\rho=yF_{yy}dy+(yF_{yb}-F_b)db\quad \Rightarrow\quad \frac{dy}{db}_{|\rho}=\frac{F_b-yF_{yb}}{yF_{yy}}
\end{equation}
Proceeding as above we get, with possible analgous rewritings in terms of the free energy density $\mathfrak{a}$,\notetag{Rewrite making manifest barotropic limit}
{\small
\begin{align}\label{eq:nonadiabatic_pressure_coefficient_gamma}
\begin{aligned}
    \Gamma:=\frac{\partial p}{\partial \sigma}_{\mid\rho}&=\frac{p_b + p_y\, \tfrac{dy}{db}_{\mid\rho}}{\sigma_b + \sigma_y\, \tfrac{dy}{db}_{\mid\rho}}=\frac{-bF_{bb}+(F_y-bF_{by})\tfrac{F_b-yF_{yb}}{yF_{yy}}}{\Big(\tfrac{bF_{by}-F_y}{b^2}\Big)+\tfrac{F_{yy}}{b}\frac{F_b-yF_{yb}}{yF_{yy}}}
    \\
    &=\frac{-b^2F_{bb}y^2F_{yy}+(yF_y-byF_{by})(bF_b-byF_{yb})}{y^2F_{yy}\Big(\tfrac{bF_{by}-F_y}{b}\Big)+yF_{yy}(F_b-yF_{yb})}
    \\
    &=\frac{b^2F_{bb}y^2F_{yy}-(yF_y-byF_{by})(bF_b-byF_{yb})}{\tfrac{1}{by}y^2F_{yy}(yF_y-bF_b)}.
\end{aligned}\end{align}}



%------------------------------------------------------------------------------
\subsubsection{Energy-momentum tensor of the fluid EFT}
%------------------------------------------------------------------------------


Note that $\Phi:=\{u^\mu, \partial^\mu \phi^i\}_{i=1,2,3}$ is a basis of the tangent spacetime.
The metric $g_{\mu\nu}$ in the coordinate basis to the metric $g^\Phi_{ab}$ in the $\Phi$-basis as
\begin{align}
     \begin{pmatrix} u^\mu \mid  \partial^\mu \phi^j\end{pmatrix}^T g_{\mu\nu} \begin{pmatrix} u^\nu  \mid  \partial^\nu \phi^j\end{pmatrix}= \begin{pmatrix} -1 & 0\\ 0 & B^{ij}\end{pmatrix} = g^{\Phi}_{ab}.
\end{align}
Inverting the matrices we find 
\begin{align}
    g^{\mu\nu}=\begin{pmatrix} u^\nu  \mid  \partial^\nu \phi^j\end{pmatrix}\begin{pmatrix} -1 & 0\\ 0 & (B^{ij})^{-1}\end{pmatrix} \begin{pmatrix} u^\mu \mid  \partial^\mu \phi^j\end{pmatrix}^T = -u^\mu u^\nu + B^{-1}_{ij} \partial^\mu \phi^i \partial^\nu \phi^j,
\end{align}
confirming that indeed $B^{-1}_{ij} \partial^\mu \phi^i \partial^\nu \phi^j= g^{\mu\nu}+u^\mu u^\nu$ is the projector on the space orthogonal to the 4-velocity.
The variation of $b=\sqrt{\det B}$ wrt $g^{\mu\nu}\mapsto g^{\mu\nu}+\delta g^{\mu\nu}$ is then
\begin{align}
    \delta b= \tfrac12 \frac{1}{\sqrt{\det B}}\, \big(\det B\big)\, \text{Tr}\Big((B^{-1})^{ij} \partial_\mu \phi^i \partial_\nu \phi^j\Big) \, \delta g^{\mu\nu} = \tfrac12 \sqrt{B}\,  (g_{\mu\nu}+u_\mu u_\nu)\, \delta g^{\mu\nu}
\end{align}
The variation of $y$ is computed similarly and gives\questiontag{Insert computation from other file}
\begin{align}
    \delta y = -\tfrac12 y \, u_\mu u_\nu \, \delta g^{\mu\nu}.
\end{align}
The energy momentum tensor is then
{\small
\begin{align}
    T_{\mu\nu}:&=-\frac{2}{\sqrt{g}}\frac{\delta \big(\sqrt{g}F\big)}{\delta g^{\mu\nu}} = -\frac{2}{\sqrt{g}}\Big[ -\tfrac12 \sqrt{g}g_{\mu\nu} F +\sqrt{g}\tfrac12 b F_b (g_{\mu\nu}+u_\mu u_\nu) - \sqrt{g}\tfrac12 y F_y u_\mu u_\nu\Big]\\[2.5pt]
    & = (yF_y - bF_b) u_\mu u_\nu + g_{\mu\nu} (F-bF_b).
\end{align}}
Notably the computation is exact to any order so that no anisotropic stress appear in thhe EFt of a single fluid.



%------------------------------------------------------------------------------
\subsubsection{Quadratic expansion of the fluid variables}
%------------------------------------------------------------------------------
The final form of the 4-velocity reads\notetag{Insert from other file}
{\small\begin{align}
\begin{aligned}
    &au^\mu=\delta_0^\mu \bigg(1 \!+\! \Big[\tfrac12 h^{00} -(\partial_j \pi^j)\Big]\! +\! \Big[ h^{0 i} \dot{\pi}^i -\tfrac12 h^{00} (\partial_j \pi^j) + \tfrac12 (\dot{\pi}^i)^2 + \tfrac12 (\partial_j \pi^j)^2 +\tfrac12 \partial_i\pi^j\partial_j\pi^i\Big]\bigg)
    \\
    &\quad\qquad+\left[1 + \tfrac12 h^{00}-(\partial_\ell\pi^\ell)\right]
     \tfrac12\epsilon_{i j k}\,\epsilon^{\mu \alpha j k}\,\partial_\alpha\pi^i 
     +\tfrac{1}{2} \epsilon_{ijk}\, \epsilon^{\mu i \beta\gamma}\partial_\beta\pi^j\,\partial_\gamma\pi^k
    + O(\pi^3, h^2)
\end{aligned}\end{align}}
This simplifies greatly if we note the identities
{\small\begin{align}\begin{aligned}
    &\epsilon_{ijk}\epsilon^{0\alpha jk}\partial_\alpha\pi^i=2\delta_i^\alpha\partial_\alpha\pi^i= 2(\partial_\ell\pi^\ell),\\[2pt]
    &\epsilon_{ijk}\, \epsilon^{0 i \beta\gamma}\partial_\beta\pi^j\,\partial_\gamma\pi^k = (\delta_j^\beta\delta_k^\gamma-\delta_j^\gamma\delta_k^\beta)\partial_\beta\pi^j\,\partial_\gamma\pi^k= (\partial_\ell\pi^\ell)^2- \partial_k\pi^j\partial_j\pi^k.
\end{aligned}
\end{align}}
The first component is then
{\small\begin{align}
\begin{aligned}
    &au^0=\bigg(1 \!+\! \Big[\tfrac12 h^{00} \bluecancel{-(\partial_j \pi^j)}\Big]\! +\! \Big[ h^{0 i} \dot{\pi}^i \redcancel{-\tfrac12 h^{00} (\partial_j \pi^j)} + \tfrac12 (\dot{\pi}^i)^2 + \greencancel{\tfrac12 (\partial_j \pi^j)^2}\graycancel{+\tfrac12 \partial_i\pi^j\partial_j\pi^i}\Big]\bigg) 
    \\
    &\bluecancel{+(\partial_\ell\pi^\ell)} \redcancel{+\tfrac12 h^{00} (\partial_\ell\pi^\ell)}\greencancel{-(\partial_\ell\pi^\ell)^2} + \greencancel{\tfrac12 (\partial_\ell\pi^\ell)^2} \graycancel{-\tfrac12 \partial_k\pi^j\partial_j\pi^k}
    \\
    &=1 + \tfrac12 h^{00} +h^{0i}\dot{\pi}^i+ \tfrac12 (\dot{\pi}^i)^2 + O(\pi^3, h^2).
\end{aligned}\end{align}}
Similarly using the identity
{\small\begin{align}\begin{aligned}
    &\epsilon_{ijk}\epsilon^{\ell\alpha jk}\partial_\alpha\pi^i= - \epsilon_{ijk}\epsilon^{0\ell jk}\delta^\alpha_0\partial_\alpha\pi^i = -2 \delta_i^\ell\delta^\alpha_0\partial_\alpha\pi^i = -2 \dot{\pi}^\ell,\\[2pt]
    & \epsilon_{ijk}\, \epsilon^{\ell i \beta\gamma}\partial_\beta\pi^j\,\partial_\gamma\pi^k= -2\delta^\beta_0 \epsilon_{ijk}\, \epsilon^{0 i \ell m} \partial_\beta\pi^j\,\partial_m\pi^k = -2 \dot{\pi}^\ell (\partial_m\pi^m) + 2 \dot{\pi}^m \partial_m\pi^\ell.
\end{aligned}
\end{align}}
the spatial components of the velocity read
{\small\begin{align}
    au^\ell &= -\left[1 + \tfrac12 h^{00}\bluecancel{-(\partial_m\pi^m)}\right]\dot{\pi}^\ell  \bluecancel{- \dot{\pi}^\ell (\partial_m\pi^m)} + \dot{\pi}^m \partial_m\pi^\ell\\
    &= - \dot{\pi}^\ell + \dot{\pi}^m \partial_m\pi^\ell - \tfrac12 h^{00}\dot{\pi}^\ell + O(\pi^3, h^2).
\end{align}}
The final form of the 4-velocity is then 
{\small\begin{align}
\begin{aligned}
    &au^\mu=\Big(1 + \tfrac12 h^{00} +h^{0i}\dot{\pi}^i+ \tfrac12 (\dot{\pi}^i)^2,\,\,- \dot{\pi}^\ell + \dot{\pi}^m \partial_m\pi^\ell - \tfrac12 h^{00}\dot{\pi}^\ell\Big)+ O(\pi^3, h^2).
\end{aligned}\end{align}}



%------------------------------------------------------------------------------
\subsubsection{Quadratic expansion of the action}
%------------------------------------------------------------------------------
For the moment the computation is reported with full details in Section \ref{sec:qudratic_action_curved_spacetime} while we figure out the most convenient form.\notetag{Eventually move it here}